\section{Formal proofs and numerical checks}\label{app:checks}

\subsection{The Lean development}

The finite inequalities of the paper are proved in Lean~4 with Mathlib, in eleven files and thirty-nine theorems, listed in Table~\ref{tab:lean}. Every theorem depends only on the three standard axioms of Lean, propositional extensionality, the axiom of choice and the soundness of quotients; the command \texttt{\#print axioms} is run on each of them, and no theorem uses \texttt{sorry}. The statements are over the real numbers, with finite sums indexed by a \texttt{Finset} and infinite ones expressed by \texttt{Summable}; in them $\sigma$ is the polynomial order written $s$ in the text, and $(1/D)^{r}$ is $\eps^{r}$ with $\eps=1/D$.

The formal statements cover the deterministic content of the main theorems: Gibbs' inequality in the form used for the entropy law, the lower bound and the two properties of the ceiling schedule (Theorem~\ref{thm:A}), the correspondence between schedules and Kraft vectors and the shifted cost (Theorem~\ref{thm:C}), H\"older's lower bound and the optimal allocation of the R\'enyi law (Theorem~\ref{thm:D}) with the summability criterion for polynomial checks, both halves of the duality criterion (Theorem~\ref{thm:dual}), the inequalities for clean runs and random schedules (Propositions~\ref{prop:clean} and~\ref{prop:random}), the precision and cost identities of the gating theorem (Theorem~\ref{thm:gating}), the summability criterion for sound singular chains with both directions (Theorem~\ref{thm:E}(iii)), the two Borel--Cantelli statements for false acceptances (Theorem~\ref{thm:E}(i)), and the divergence and amplification statements of the branching theorem (Theorem~\ref{thm:F}). The limit theorems of probability (Wald's identity, the extinction theorem, the Kesten--Stigum theorem and L\'evy's lemma), the asymptotics of $\zeta$ and of the entropies in Propositions~\ref{prop:powerlaw} and~\ref{prop:osgood}, and the estimates of Proposition~\ref{prop:twosided} are proved in the text and are not formalised.

\begin{table*}[t]
\caption{The thirty-nine theorems of the Lean development, by file. The third column gives the statement in the notation of the paper and the fourth the result in which it is used. Sums over $i$ are finite; sums over $n$ or $k$ are infinite and ``$\sum<\infty$'' means summability.}
\label{tab:lean}
\centering
\footnotesize
\begin{tabular}{@{}l l p{0.53\textwidth} l@{}}
\toprule
File & Theorem & Statement & Used in \\
\midrule
\texttt{Basic} & \lean{gibbs\_term} & $x\log\frac1x\le x\log\frac1y+y-x$ for $x,y>0$ & Lemma~\ref{lem:gibbs} \\
 & \lean{entropy\_term} & $a\log\frac Aa+a\log\frac1\theta\le a\log\frac1b+\frac{Ab}\theta-a$ & Lemma~\ref{lem:entropyterm} \\
 & \lean{log\_inv\_pow} & $\log(1/\eps^{r})=r\log(1/\eps)$ & Theorem~\ref{thm:A} \\
 & \lean{log\_inv\_pos} & $\log(1/\eps)>0$ for $0<\eps<1$ & Theorem~\ref{thm:A} \\
\texttt{EntropyLaw} & \lean{entropy\_lower} & $\sum\beta_{i}\le\theta\Rightarrow\sum a_{i}\log\frac A{a_{i}}+A\log\frac1\theta\le\sum a_{i}\log\frac1{\beta_{i}}$ & Proposition~\ref{prop:entropyineq} \\
 & \lean{verification\_lower} & $\sum\eps^{r_{i}}\le\theta\Rightarrow\frac AL(H+\log\frac1\theta)\le\sum a_{i}r_{i}$ & Theorem~\ref{thm:A} \\
 & \lean{ceilSchedule\_sound} & the ceiling schedule has $\sum\eps^{r_{i}}\le\theta$ & Theorem~\ref{thm:A} \\
 & \lean{ceilSchedule\_cost} & the ceiling schedule costs at most $\frac AL(H+\log\frac1\theta)+A$ for $\theta\le1$ & Theorem~\ref{thm:A} \\
\texttt{Kraft} & \lean{kraft\_of\_schedule} & $\sum D^{-r_{i}}\le D^{-k}\Rightarrow r_{i}\ge k$ and $\sum D^{-(r_{i}-k)}\le1$ & Lemma~\ref{lem:kraftsound}(i) \\
 & \lean{schedule\_of\_kraft} & $\sum D^{-\ell_{i}}\le1\Rightarrow\sum D^{-(\ell_{i}+k)}\le D^{-k}$ & Lemma~\ref{lem:kraftsound}(ii) \\
 & \lean{shifted\_cost} & $\sum a_{i}(\ell_{i}+k)=kA+\sum a_{i}\ell_{i}$ & Lemma~\ref{lem:kraftsound}(iii) \\
\texttt{Renyi} & \lean{renyi\_lower} & $\sum\delta_{i}\le\theta\Rightarrow(\sum a_{i}^{1/(1+s)})^{1+s}\le\theta^{s}\sum a_{i}\delta_{i}^{-s}$ & Theorem~\ref{thm:D} \\
 & \lean{renyi\_attain} & $\delta_{i}=\theta a_{i}^{1/(1+s)}/S$ has $\sum\delta_{i}=\theta$ and cost $\theta^{-s}S^{1+s}$ & Theorem~\ref{thm:D} \\
 & \lean{poly\_schedule\_iff} & a positive $\delta$ with $\sum\delta_{n}<\infty$, $\sum a_{n}\delta_{n}^{-s}<\infty$ exists $\iff\sum a_{n}^{1/(1+s)}<\infty$ & Corollary~\ref{cor:dualcases}(ii) \\
\texttt{Duality} & \lean{dualCost\_nonneg} & $\psi(u)\ge0$ for $u\ge0$ & Theorem~\ref{thm:dual} \\
 & \lean{dualCost\_le} & $\psi(u)\le\kappa(\delta)+u\delta$ for $\delta\in(0,1]$ & Theorem~\ref{thm:dual} \\
 & \lean{dual\_summable\_of\_schedule} & a sound allocation of finite cost gives $\sum a_{n}\psi(\lambda/a_{n})<\infty$ for every $\lambda>0$ & Theorem~\ref{thm:dual} \\
 & \lean{schedule\_of\_dual\_summable} & $\sum a_{n}\psi(\lambda/a_{n})<\infty$ for one $\lambda>0$ gives a sound allocation of finite cost & Theorem~\ref{thm:dual} \\
\texttt{Clean} & \lean{clean\_le\_exp} & $\prod(1-q_{i})\le e^{-\sum q_{i}}$ for $q_{i}\le1$ & Proposition~\ref{prop:clean} \\
 & \lean{clean\_ge} & $1-\sum q_{i}\le\prod(1-q_{i})$ for $0\le q_{i}\le1$ & Proposition~\ref{prop:clean} \\
 & \lean{clean\_budget} & $\prod(1-q_{i})\ge1-\theta\Rightarrow\sum q_{i}\le\log\frac1{1-\theta}$ & Proposition~\ref{prop:clean} \\
\texttt{Adaptive} & \lean{jensen\_repetition} & $\eps^{\sum w_{j}r_{j}}\le\sum w_{j}\eps^{r_{j}}$ for weights $w$ & Proposition~\ref{prop:random} \\
 & \lean{randomised\_budget} & a random schedule within budget in mean has its mean within budget & Proposition~\ref{prop:random} \\
\texttt{Gating} & \lean{step\_precision} & $\frac{(1-\rho)(1-\eta)}{(1-\rho)(1-\eta)+\rho\delta}=\frac1{1+x}$ & Theorem~\ref{thm:gating}(i) \\
 & \lean{precision\_bounds} & $e^{-\sum x_{i}}\le\prod(1+x_{i})^{-1}\le e^{-\sum x_{i}/(1+x_{i})}$ & Theorem~\ref{thm:gating}(i) \\
 & \lean{outcome\_precision} & $\frac{c(1-\eta_{o})}{c(1-\eta_{o})+(1-c)\delta_{o}}=\bigl(1+\frac{(1-c)\delta_{o}}{c(1-\eta_{o})}\bigr)^{-1}$ & Theorem~\ref{thm:gating}(ii) \\
 & \lean{separation} & cost ratio $(1-\rho)^{-n}$ and $1+n\rho\le(1-\rho)^{-n}$ & Theorem~\ref{thm:gating}(iii) \\
\texttt{Singularity} & \lean{entropy\_summable\_of\_schedule} & $\sum\eps^{r_{n}}<\infty$, $\sum a_{n}r_{n}<\infty$, $a_{n}\le1\Rightarrow\sum a_{n}\log\frac1{a_{n}}<\infty$ & Theorem~\ref{thm:E}(iii) \\
 & \lean{costSchedule\_pow\_le} & $r_{n}=\lceil\log(1/a_{n})/L\rceil$ has $\eps^{r_{n}}\le a_{n}$ & Proposition~\ref{prop:integer} \\
 & \lean{schedule\_of\_entropy\_summable} & $\sum a_{n}<\infty$, $\sum a_{n}\log\frac1{a_{n}}<\infty$ make that schedule sound and of finite cost & Theorem~\ref{thm:E}(iii) \\
 & \lean{sound\_schedule\_iff} & for $\sum a_{n}<\infty$, $a_{n}\le1$: a sound schedule of finite cost exists $\iff\sum a_{n}\log\frac1{a_{n}}<\infty$ & Proposition~\ref{prop:integer} \\
 & \lean{shift\_schedule} & $\eps^{r_{n}+k}=\eps^{k}\eps^{r_{n}}$ & Theorem~\ref{thm:E}(iii) \\
\texttt{Soundness} & \lean{false\_accepts\_finite} & $\PR(E_{n})\le Cq_{n}$, $\sum q_{n}<\infty\Rightarrow\PR(E_{n}\text{ i.o.})=0$ & Theorem~\ref{thm:E}(i) \\
 & \lean{false\_accepts\_infinite} & independent $E_{n}$, $\PR(E_{n})\ge bq_{n}$, $b>0$, $\sum q_{n}=\infty\Rightarrow\PR(E_{n}\text{ i.o.})=1$ & Theorem~\ref{thm:E}(i) \\
 & \lean{any\_false\_accept\_le} & $\PR(\bigcup E_{n})\le C\sum q_{n}$ & Theorem~\ref{thm:E}(i) \\
\texttt{Branching} & \lean{contamination\_diverges} & $\sum m^{k}\delta=\infty$ for $m\ge1$, $\delta>0$ & Theorem~\ref{thm:F} \\
 & \lean{amplified\_term} & $k\log(2m)\le r_{k}L\Rightarrow m^{k}\eps^{r_{k}}\le2^{-k}$ & Theorem~\ref{thm:F} \\
 & \lean{amplified\_summable} & under the same hypothesis $\sum m^{k}\eps^{r_{k}}<\infty$ & Theorem~\ref{thm:F} \\
 & \lean{amplified\_schedule\_ok} & $r_{k}=\lceil k\log(2m)/L\rceil$ satisfies that hypothesis & Theorem~\ref{thm:F} \\
\bottomrule
\end{tabular}
\end{table*}

\subsection{Numerical checks}

Independent programs test the closed forms of the paper against exhaustive searches, exact optima and simulation, on instances drawn with fixed seeds. They are listed in Table~\ref{tab:checks}; the shell script that runs them also builds the Lean development and prints the axioms of every theorem, and it reports $33{,}505$ checks with no failure. The checks are not part of any proof. They guard against errors in the statements: a wrong constant, a reversed inequality or a misprinted formula would make one of them fail. The numerical values quoted in the text, in Tables~\ref{tab:huffman} and~\ref{tab:gw} and in the figures are produced by the same programs.

\begin{table}[!htb]
\caption{The numerical checks.}
\label{tab:checks}
\centering
\footnotesize
\begin{tabular}{@{}l >{\raggedright\arraybackslash}p{0.56\columnwidth} r@{}}
\toprule
Language & What is checked & Checks \\
\midrule
Python & Theorems~\ref{thm:A}, \ref{thm:C}: bounds and Huffman cost against exact optima & 649 \\
Python & Theorem~\ref{thm:gating}: costs and precisions against simulation & 533 \\
Python & Theorems~\ref{thm:D}, \ref{thm:dual}: R\'enyi optimum, dual costs & 20{,}850 \\
Python & Theorem~\ref{thm:E}, Propositions~\ref{prop:powerlaw}, \ref{prop:osgood} & 241 \\
Python & Theorem~\ref{thm:F}: takeoff, kept errors, amplified checks & 2{,}432 \\
C & Theorem~\ref{thm:C}: exhaustive search over schedules & 800 \\
C\texttt{++} & Theorems~\ref{thm:D}, \ref{thm:gating} at random feasible points & 8{,}000 \\
\midrule
 & total & 33{,}505 \\
\bottomrule
\end{tabular}
\end{table}
